% Template for PLoS
% Version 3.8 Apr 2026
%
% % % % % % % % % % % % % % % % % % % % % %
%
% -- IMPORTANT NOTE
%
% This template contains comments intended 
% to minimize problems and delays during our production 
% process. Please follow the template instructions
% whenever possible.
%
% % % % % % % % % % % % % % % % % % % % % % % 
%
% Once your paper is accepted for publication, 
% PLEASE REMOVE ALL TRACKED CHANGES in this file 
% and leave only the final text of your manuscript. 
% PLOS recommends the use of latexdiff to track changes during review, as this will help to maintain a clean tex file.
% Visit https://www.ctan.org/pkg/latexdiff?lang=en for info.
%
%
% IMPORTANT
% Below are a few tips to help format manuscript according to our specifications. For more tips to help reduce the possibility of formatting errors during conversion, please see our LaTeX guidelines at http://journals.plos.org/plosone/s/latex
%
% There are no restrictions on package use within the LaTeX files except that no packages listed in the template may be deleted.
%
% Please do not include colors or graphics in the text.
%
% The manuscript LaTeX source should be contained within a single file (do not use \input, \externaldocument, or similar commands).
%
% % % % % % % % % % % % % % % % % % % % % % %
%
% -- FIGURES AND TABLES
%
% Please include tables/figure captions directly after the paragraph where they are first cited in the text.
%
% DO NOT INCLUDE GRAPHICS IN YOUR MANUSCRIPT
% - Figures should be uploaded separately from your manuscript file. 
% - Figures generated using LaTeX should be extracted and removed from the PDF before submission. 
% - Figures containing multiple panels/subfigures must be combined into one image file before submission.
% For figure citations, please use "Fig" instead of "Figure".
% See http://journals.plos.org/plosone/s/figures for PLOS figure guidelines.
%
% Tables should be cell-based and may not contain:
% - spacing/line breaks within cells to alter layout or alignment
% - do not nest tabular environments (no tabular environments within tabular environments)
% - no graphics or colored text (cell background color/shading OK)
% See http://journals.plos.org/plosone/s/tables for table guidelines.
%
% For tables that exceed the width of the text column, use the adjustwidth environment as illustrated in the example table in text below.
%
% % % % % % % % % % % % % % % % % % % % % % % %
%
% -- EQUATIONS, MATH SYMBOLS, OTHER SPECIAL FORMATTING
%
% - For bold symbols, use the \boldsymbol command and not \mathbf. 
% - For inline equations, please be sure to include all portions of an equation in the math environment.  For example, x$^2$ is incorrect; this should be formatted as $x^2$ (or $\mathrm{x}^2$ if the romanized font is desired).
% - Do not include text that is not math in the math environment. For example, CO2 should be written as CO\textsubscript{2} instead of CO$_2$.
% - Avoid capturing simple text as equations, for example $<$.
% - For inline equations, please do not include punctuation (commas, etc) within the math environment unless this is part of the equation.
% - When adding superscript or subscripts outside of brackets/braces, please group using {}.  For example, change "[U(D,E,\gamma)]^2" to "{[U(D,E,\gamma)]}^2". 
% - Do not use \cal for caligraphic font.  Instead, use \mathcal{}.
% - Avoid splitting for fragmenting a single equation into multiple objects/environments, for example $f(x)$ = $g(x)$ + $h(x)$.
% - Use math environment for all equations, do not mimic using text.
% - Avoid using single letters in math mode where possible.
% - Always use matching parentheses and delimiters, avoid mixing math and text brackets.
% - Insert spaces around inline equations to ensure proper display after conversion.
% - Do not use forced numbers for display equation labeling or suppress numbering with \nonumber. 
%
% % % % % % % % % % % % % % % % % % % % % % % % 
%
% Please contact customercare@plos.org with any questions.
%
% % % % % % % % % % % % % % % % % % % % % % % %

\documentclass[10pt,letterpaper]{article}
\usepackage[top=0.85in,left=2.75in,footskip=0.75in]{geometry}

% amsmath and amssymb packages, useful for mathematical formulas and symbols
\usepackage{amsmath,amssymb}

% Use adjustwidth environment to exceed column width (see example table in text)
\usepackage{changepage}

% textcomp package and marvosym package for additional characters
\usepackage{textcomp,marvosym}

% cite package, to clean up citations in the main text. Do not remove.
\usepackage{cite}

% Use nameref to cite supporting information files (see Supporting Information section for more info)
\usepackage{nameref,hyperref}

% line numbers
\usepackage[right]{lineno}

% ligatures disabled
\usepackage[nopatch=eqnum]{microtype}
\DisableLigatures[f]{encoding = *, family = * }

% color can be used to apply background shading to table cells only
\usepackage[table]{xcolor}

% array package and thick rules for tables
\usepackage{array}

% user packages
\usepackage{booktabs}
\usepackage{pdflscape}
% Define the colors used in the CPE confusion table
\definecolor{TPgreen}{HTML}{98FB98}
\definecolor{TNwhite}{HTML}{FFFFFF}
\definecolor{FNyellow}{HTML}{FFFACD}
\definecolor{FPred}{HTML}{FFB6C1}

% create "+" rule type for thick vertical lines
\newcolumntype{+}{!{\vrule width 2pt}}

% create \thickcline for thick horizontal lines of variable length
\newlength\savedwidth
\newcommand\thickcline[1]{%
  \noalign{\global\savedwidth\arrayrulewidth\global\arrayrulewidth 2pt}%
  \cline{#1}%
  \noalign{\vskip\arrayrulewidth}%
  \noalign{\global\arrayrulewidth\savedwidth}%
}

% \thickhline command for thick horizontal lines that span the table
\newcommand\thickhline{\noalign{\global\savedwidth\arrayrulewidth\global\arrayrulewidth 2pt}%
\hline
\noalign{\global\arrayrulewidth\savedwidth}}


% Remove comment for double spacing
%\usepackage{setspace} 
%\doublespacing

% Text layout
\raggedright
\setlength{\parindent}{0.5cm}
\textwidth 5.25in
\textheight 8.75in

% Bold the 'Figure #' in the caption and separate it from the title/caption with a period
% Captions will be left justified
\usepackage[aboveskip=1pt,labelfont=bf,labelsep=period,justification=raggedright,singlelinecheck=off]{caption}
\renewcommand{\figurename}{Fig}

% Please use the included `plos2025.bst` as your BibTeX style
\bibliographystyle{plos2025}

% Remove brackets from numbering in List of References
\makeatletter
\renewcommand{\@biblabel}[1]{\quad#1.}
\makeatother



% Header and Footer with logo
\usepackage{lastpage,fancyhdr,graphicx}
\usepackage{epstopdf}
%\pagestyle{myheadings}
\pagestyle{fancy}
\fancyhf{}
%\setlength{\headheight}{27.023pt}
%\lhead{\includegraphics[width=2.0in]{PLOS-submission.eps}}
\rfoot{\thepage/\pageref{LastPage}}
\renewcommand{\headrulewidth}{0pt}
\renewcommand{\footrule}{\hrule height 2pt \vspace{2mm}}
\fancyheadoffset[L]{2.25in}
\fancyfootoffset[L]{2.25in}
\lfoot{\today}

%% Include all user-defined macros below

\newcommand{\lorem}{\textbf{LOREM}}
\newcommand{\ipsum}{\textbf{IPSUM}}

%% END MACROS SECTION


\begin{document}
\vspace*{0.2in}

% Title must be 250 characters or less.
\begin{flushleft}
{\Large
\textbf{Vero cell culture CPE detection benchmarking} % Please use "sentence case" for title and headings (capitalize only the first word in a title (or heading), the first word in a subtitle (or subheading), and any proper nouns).
}
\newline
% Insert author names, affiliations and corresponding author email (do not include titles, positions, or degrees).
\\
Albert Mathews\textsuperscript{1}
\\
\bigskip
\textbf{1} Independent Researcher, Montreal, Quebec, Canada
\\
\bigskip

% Insert additional author notes using the symbols described below. Insert symbol callouts after author names as necessary.
% 
% Remove or comment out the symbols in the byline and the author notes below if they aren't used.
%
% Primary Equal Contribution Note
%\Yinyang These authors contributed equally to this work.

% Additional Equal Contribution Note
% Also use this double-dagger symbol for special authorship notes, such as senior authorship.
%\ddag These authors also contributed equally to this work.

% Current address notes
%\textcurrency Current Address: Dept/Program/Center, Institution Name, City, State, Country % change symbol to "\textcurrency a" if more than one current address note
% \textcurrency b Insert second current address 
% \textcurrency c Insert third current address

% Deceased author note
%\dag Deceased

% Group/Consortium Author Note
%\textpilcrow Membership list can be found in the Acknowledgments section.

% Use the asterisk to denote corresponding authorship and provide email address in note below.
* albert.mathews@mail.mcgill.ca

\end{flushleft}
% Please keep the abstract below 300 words

% For PLOS Medicine research article authors, please structure your abstract
% with "Background", "Method and Findings" and "Conclusion" sections per
% journal requirements.

% For PLOS Neglected Tropical Diseases research article authors, please
% structure your abstract with "Background", "Methodology", "Findings", and
% "Conclusion" sections per journal requirements.
%
\section*{Abstract}
This study evaluates the performance of three domain-specific computer-vision tools (AIRVIC, DVICE, and Cellpose) and four general-purpose multimodal large language models (ChatGPT, Claude, Gemini, and Grok) on the task of detecting and classifying cytopathic effect (CPE) in Vero-cell culture microscope images. Ground-truth CPE presence and morphological types were derived solely from narrative descriptions supplied by the contract research organization (CRO) that performed the cell-culture experiments and were cross-referenced against the standardized CPE descriptors listed in Table 7 of the CLSI M41-A Guideline. Quantitative analysis was performed exclusively on the 22 images that received CRO descriptions.

None of the evaluated tools demonstrated adequate performance for reliable CPE detection. Domain-specific tools showed extremely high false-positive rates, while the multimodal models exhibited either strong conservative bias or only modest balanced accuracy. These findings indicate that current AI approaches, whether domain-specific or general-purpose, are not suitable as objective tools for automated CPE analysis of the Vero cell culture image dataset used in this study.


\clearpage
\newgeometry{top=0.85in,left=1in,right=1in,footskip=0.75in}
\linenumbers

\section*{Introduction}
\label{sec:intro}
This study benchmarks the performance of three domain-specific computer-vision tools (AIRVIC, DVICE, and Cellpose) and four general-purpose multimodal large language models (ChatGPT, Claude, Gemini, and Grok) on the task of detecting and classifying cytopathic effect (CPE) in Vero-cell culture microscope images. The objective is to evaluate whether any of these AI approaches can serve as objective, reproducible analysis tools for CPE detection in previously unseen images. Ground-truth CPE status and morphological types were derived exclusively from the image descriptions supplied by the contract research organization (CRO) that performed the cell-culture experiments. These descriptions were cross-referenced against the standardized CPE descriptors listed in Table 7 of the CLSI M41-A Guideline.
\section*{Materials and methods}
\label{sec:methods}
All code and scripts are available in the project GitHub repository \cite{github-verocell}.

\subsection*{Image dataset and ground truth}
The study utilized the publicly available Vero Cell Culture Image Dataset \cite{mathews2025}. This dataset consists of images acquired in two distinct stages: a preparation (PREP) stage in which cells were thawed and passaged until normal growth rates were achieved, followed by an experimental (EXP) stage. The present study focuses exclusively on the 101 images from the EXP stage. These images were obtained from two experimental paths (Path 1 maintained at 10\% FBS and Path 2 reduced to 2\% FBS) initiated at passage 4 (see sample images \nameref{S1_Fig} and \nameref{S2_Fig}). The CRO was instructed to provide descriptive captions for any image that would benefit from accompanying text; no specific directive was given to search for, classify, or quantify cytopathic effects (CPE). Consequently, only 22 of the 101 EXP-stage images contain CRO descriptions (hereafter referred to as the ``description set''), while the remaining 83 do not. Accuracy analyses were performed exclusively on the description set.

\subsection*{Standard CPE descriptors}
Characteristic CPE descriptors were extracted from Table 7 of the CLSI M41-A Guideline \cite{clsi} through an iterative review process. First, the full table was examined to identify a set of ten descriptors: rounded, ballooned/enlarged, syncytia, vacuolation, detachment, granularity, refractile, cytoplasmic strands, nonspecific degeneration, and ``no CPE''. An incidence table was then constructed by identifying every instance of these terms (or close semantic correlates) in the ``Appearance'' column (see Table \ref{tab:clsi-m41-cpe}).

\begin{landscape}
\begin{table}[htbp]
\centering
\small
\caption{CLSI M41 Table 7. Characteristic CPE in [Inoculated] Tube Cultures}
\label{tab:clsi-m41-cpe}

\begin{tabular}{l *{10}{>{\centering\arraybackslash}m{0.5cm}} p{3.2cm}}
\toprule
\textbf{Virus} 
& \rotatebox{90}{\textbf{Rounded}} 
& \rotatebox{90}{\textbf{Ballooned/Enlarged}} 
& \rotatebox{90}{\textbf{Syncytia}} 
& \rotatebox{90}{\textbf{Vacuolation}} 
& \rotatebox{90}{\textbf{Detached}} 
& \rotatebox{90}{\textbf{Granularity}} 
& \rotatebox{90}{\textbf{Refractile}} 
& \rotatebox{90}{\textbf{Cytoplasmic strands}} 
& \rotatebox{90}{\textbf{Nonspecific degeneration}} 
& \rotatebox{90}{\textbf{No CPE**}} \\

\midrule
Adenovirus          & x & x &   &   &   &   &   &   &   &   \\
Cytomegalovirus     & x & x &   &   &   &   &   &   &   &   \\
Enterovirus         & x &   &   &   &   &   & x &   &   &   \\
Herpes Simplex      & x & x & x &   &   &   &   &   &   &   \\
Influenza           &   &   &   & x &   & x &   &   & x & x \\
Measles             &   & x & x & x &   & x &   &   &   &   \\
Metapneumovirus     & x &   & x &   & x &   & x &   &   &   \\
Mumps               & x &   & x &   &   & x &   &   & x &   \\
Parainfluenza viruses & x &   & x &   &   & x &   &   & x &   \\
Reoviruses          &   &   &   &   & x &   &   &   & x &   \\
Respiratory Syncytial Virus &   &   & x &   &   & x &   &   & x &   \\
Rhinoviruses        & x &   &   &   &   &   & x &   &   &   \\
SARS-CoV            & x &   &   &   & x &   & x &   &   &   \\
Vaccinia            & x & x & x &   &   &   & x & x &   &   \\
Varicella-zoster virus & x & x & x &   &   & x & x & x &   &   \\
\bottomrule
\end{tabular}

\end{table}
\end{landscape}
Next, the same mapping procedure was applied to the CRO image descriptions to generate a second incidence table. Only five of the original ten CLSI descriptors appeared in the CRO descriptions. A sixth term, ``Dying Cells'', emerged directly from the CRO text and was retained as a distinct category (Dy) rather than being forced into the CLSI ``nonspecific degeneration'' bin. The resulting \textit{evaluation descriptor set} used for ground truth and all subsequent analyses therefore consisted of: Dying cells (Dy), Rounded (Ro), Vacuolation (V), Detached (D), Granularity (G), and Refractile (Re). The CRO incidence table \ref{tab:cro-cpe-detections} below shows the result of applying this evaluation descriptor set to the image descriptions.

\begin{table}[htbp]
\centering
\small
\caption{CRO-derived ground-truth CPE detections from image descriptions (evaluation descriptor set). ``x'' indicates presence of the corresponding CPE morphological descriptor; blank cells indicate absence.}
\label{tab:cro-cpe-detections}

\begin{tabular}{cc*{6}{c}}
\toprule
\textbf{Path} & \textbf{ID} & \textbf{Dy} & \textbf{Ro} & \textbf{V} & \textbf{D} & \textbf{G} & \textbf{Re} \\
\midrule
1 & 101 &  &  &  &  &  &  \\
1 & 103 &  &  &  &  &  &  \\
1 & 104 &  &  &  &  &  &  \\
1 & 201 &  &  &  &  &  &  \\
1 & 202 &  &  &  &  &  &  \\
1 & 305 &  &  &  &  &  &  \\
1 & 401 &  &  &  &  &  &  \\
1 & 406 &  &  & x &  &  &  \\
1 & 501 &  &  &  &  &  &  \\
2 & 101 &  & x &  &  &  &  \\
2 & 201 &  &  &  &  &  &  \\
2 & 202 &  & x &  &  &  & x \\
2 & 303 &  & x & x &  &  &  \\
2 & 307 &  &  &  &  &  &  \\
2 & 308 &  &  &  &  &  &  \\
2 & 402 & x & x & x &  &  & x \\
2 & 405 & x & x &  &  & x &  \\
2 & 406 & x & x &  &  & x &  \\
2 & 503 & x &  & x & x & x & x \\
2 & 504 & x &  & x & x & x & x \\
2 & 507 & x & x &  &  & x &  \\
2 & 508 & x & x &  &  & x &  \\
\bottomrule
\end{tabular}
\end{table}

\subsection*{AIRVIC analysis}
AIRVIC \cite{airvic} is a ResNet-based classifier trained specifically for CPE detection and virus identification in Vero and MDBK cell-culture images. All 101 images were uploaded to the web interface \cite{airvicweb} with \textit{cell line type} set to ``Vero'' and \textit{virus type} set to ``Unknown''. Because the interface does not support batch export, results were extracted manually and stored in CSV format. AIRVIC outputs a binary CPE decision (present/absent) and a predicted virus identity; the virus predictions are not used in this study, which focuses exclusively on CPE detection.

\subsection*{DVICE analysis}
DVICE \cite{dvice} is a versatile automated pipeline for label-free virus infectivity quantification using deep learning on light-microscopy images. Access to three pre-trained models was provided by the authors. All 101 EXP-stage images were processed by each of the three models. The final DVICE CPE probability score was computed as the average of the three model outputs. A threshold sweep was performed on the description set to evaluate the sensitivity of accuracy to the decision threshold. Binary CPE detection was determined using the conventional threshold of 0.5 applied to the averaged probability. The exact code, model outputs, and full threshold sweep are available in the project GitHub repository \cite{github-verocell}.

\subsection*{Cellpose analysis}
Cellpose \cite{cellpose} was used as the primary segmentation engine. A custom Python script, developed iteratively in collaboration with Grok \cite{grok}, processed all 101 images to generate instance masks---binary images in which each individual cell is delineated as a separate labeled region. Post-processing metrics were then computed on the masks to derive a proxy for CPE presence. Circularity \((4\pi \cdot \text{area} / \text{perimeter}^2)\) was selected because it is a standard morphometric parameter that effectively quantifies the transition from spread polygonal to rounded cell shape, a morphological change widely recognized as a hallmark of cytopathic effect (CPE) in cell cultures \cite{clsi} \cite{fukaya}. Eccentricity was retained as a complementary descriptor because prior label-free CPE analyses in Vero cultures have identified cell-rounding morphology as a robust, objective visual proxy for cytopathic changes, which can be quantified using shape descriptors such as aspect ratio \cite{revvity}. These two metrics were therefore chosen as the most suitable label-free proxies supported by the literature.

\begin{table}[htbp]
\centering
\small
\setlength{\tabcolsep}{8pt}
\caption{CPE proxy metrics used for CellPose probability and detection (circularity and eccentricity only)}
\label{tab:cpe-proxies-final}

\begin{tabular}{@{}p{5.2cm}p{2.8cm}p{2.8cm}p{5.5cm}@{}}
\toprule
\textbf{Description} 
& \textbf{Healthy Range (10\% FBS)} 
& \textbf{CPE Range (2\% FBS)} 
& \textbf{Rationale as CPE Proxy} \\
\midrule
Degree of cell roundness (4π·area/perimeter², 0--1) 
& 0.3--0.6 
& 0.7--1.0 
& Higher values directly quantify the transition from spread polygonal to rounded morphology, a core visual CPE feature; independent of FBS-driven growth-rate changes. \\
Shape elongation descriptor (0 = circle, 1 = line) 
& 0.4--0.7 
& 0.8--0.95 
& Captures loss of elongation; increased eccentricity indicates rounding and is a robust, label-free proxy for cytopathic changes; unaffected by the intentional 2\% FBS reduction in Path 2. \\
\bottomrule
\end{tabular}
\end{table}

The circularity and eccentricity metrics were post-processed to derive a continuous CPE probability score. Binary CPE detection was obtained using the conventional threshold of 0.5 (values $>$ 0.5 classified as positive). A threshold sweep was performed across the description set to evaluate sensitivity to the decision threshold. The exact probability formula, post-processing code, and full threshold sweep are available in the project GitHub repository \cite{github-verocell}.

\subsection*{Multimodal AI analysis}
Four multimodal large language models \cite{chatgpt} \cite{claude} \cite{gemini} \cite{grok} were evaluated under identical conditions. For each model a dedicated Python script was written that (1) loads the PNG images, (2) supplies a standardized prompt requesting detection of the six \textit{evaluation descriptor set} morphologies, and (3) requests formatted output. Each model was offered two example images with known ground truth for optional in-context learning. Prompt and processing-script engineering was performed collaboratively with the respective model (i.e. models' are co-authors of the script and prompt). Model outputs were post-processed by a script to produce per-image, per-morphology-type binary decisions that were compared against CRO ground truth.

\subsection*{Other tools considered and excluded}
Standard image-processing tools commonly used in cell-culture analysis were evaluated but excluded for the following reasons. Fiji/ImageJ \cite{fiji} requires manual filter tuning or macro development that would necessarily incorporate knowledge of the CRO descriptions, introducing circular reasoning and violating the requirement for an objective, a-priori tool. CellProfiler \cite{cellprofiler} was tested with the ``ExampleHuman'' pipeline; however, the pipeline proved too rigid for the present image set, and any custom pipeline development would again risk circular reasoning. Snapcyte \cite{snapcyte} could not be accessed after repeated failed account-creation attempts and lack of response from the vendor. Revvity Signals Image Artist \cite{revvityIA} was discussed with the vendor but commercial and technical setup barriers prevented its inclusion. These exclusions are reported for transparency; the tools ultimately benchmarked represent the most readily available and least biased options that could be applied without post-hoc tuning to the ground truth.

\subsection*{Statistical evaluation}
Performance was quantified using false-positive rate, false-negative rate, true-positive rate, true-negative rate, and overall accuracy. For the multimodal models these metrics were first computed separately for each of the six CPE morphological types and then macro-averaged. For AIRVIC, DVICE, and Cellpose the metrics were calculated directly on the binary CPE detection task. Baseline ``always-positive'' and ``always-negative'' classifiers were also evaluated for comparison. All calculations were performed on the 22-image \textit{description set} only.
\section*{Results}
\label{sec:results}
Performance results for each tool are presented below.
\clearpage
\subsection*{AIRVIC}
The CPE detection and virus prediction results obtained with AIRVIC are shown in Table~\ref{tab:airvic-results}.
\begin{table}[htbp]
\centering
\small
\setlength{\tabcolsep}{8pt}
\caption{AIRVIC CPE detection and virus prediction results.}
\label{tab:airvic-results}

\begin{tabular}{ccccc}
\toprule
\multicolumn{2}{c}{\textbf{Image}} 
& \textbf{CRO} 
& \multicolumn{2}{c}{\textbf{AIRVIC}} \\
\cmidrule(lr){1-2}\cmidrule(lr){3-3}\cmidrule(lr){4-5}
path & id 
& CPE 
& CPE & Virus \\
\midrule
1 & 101 & - & x & SARS-CoV-2 \\
1 & 103 & - & x & SARS-CoV-2 \\
1 & 104 & - & x & SARS-CoV-2 \\
1 & 201 & - & - & none \\
1 & 202 & - & x & SARS-CoV-2 \\
1 & 305 & - & x & SARS-CoV-2 \\
1 & 401 & - & x & SARS-CoV-2 \\
1 & 406 & x & x & SARS-CoV-2 \\
1 & 501 & - & x & BoAHV-1 \\
2 & 101 & x & x & SARS-CoV-2 \\
2 & 201 & - & x & SARS-CoV-2 \\
2 & 202 & x & x & SARS-CoV-2 \\
2 & 303 & x & x & SARS-CoV-2 \\
2 & 307 & - & x & BoAHV-1 \\
2 & 308 & - & x & SARS-CoV-2 \\
2 & 402 & x & x & SARS-CoV-2 \\
2 & 405 & x & x & BoAHV-1 \\
2 & 406 & x & x & SARS-CoV-2 \\
2 & 503 & x & x & SARS-CoV-2 \\
2 & 504 & x & x & SARS-CoV-2 \\
2 & 507 & x & x & BoAHV-1 \\
2 & 508 & x & x & SARS-CoV-2 \\
\bottomrule
\end{tabular}
\end{table}
\clearpage

\subsection*{DVICE}
The following two tables present the per-image CPE probability scores and performance metrics obtained with DVICE.
\begin{table}[htbp]
\centering
\small
\setlength{\tabcolsep}{8pt}
\caption{DVICE CPE probability results.}
\label{tab:dvice-results}

\begin{tabular}{ccccc}
\toprule
\multicolumn{2}{c}{\textbf{Image}} 
& \textbf{CRO} 
& \multicolumn{2}{c}{\textbf{DVICE}} \\
\cmidrule(lr){1-2}\cmidrule(lr){3-3}\cmidrule(lr){4-5}
path & id & CPE & Avg Probability & CPE Detection \\
\midrule
\arrayrulecolor{gray!30}
1 & 101 & - & 0.9997 & x \\
\hline
1 & 103 & - & 0.5331 & x \\
\hline
1 & 104 & - & 0.6151 & x \\
\hline
1 & 201 & - & 1.0000 & x \\
\hline
1 & 202 & - & 0.8871 & x \\
\hline
1 & 305 & - & 0.9998 & x \\
\hline
1 & 401 & - & 1.0000 & x \\
\hline
1 & 406 & x & 1.0000 & x \\
\hline
1 & 501 & - & 1.0000 & x \\
\hline
2 & 101 & x & 0.6668 & x \\
\hline
2 & 201 & - & 0.6666 & x \\
\hline
2 & 202 & x & 0.6326 & x \\
\hline
2 & 303 & x & 0.4452 & - \\
\hline
2 & 307 & - & 0.3369 & - \\
\hline
2 & 308 & - & 0.2723 & - \\
\hline
2 & 402 & x & 0.7760 & x \\
\hline
2 & 405 & x & 0.6612 & x \\
\hline
2 & 406 & x & 0.6638 & x \\
\hline
2 & 503 & x & 0.8841 & x \\
\hline
2 & 504 & x & 0.7337 & x \\
\hline
2 & 507 & x & 0.6543 & x \\
\hline
2 & 508 & x & 0.6635 & x \\
\arrayrulecolor{black}
\bottomrule
\end{tabular}

\end{table}
\begin{table}[htbp]
\centering
\small
\caption{DVICE CPE detection performance metrics (description set, n=22). Three models were provided by the authors; the average probability across the three models was used for scoring and final binary decision.}
\label{tab:dvice-performance}

\begin{tabular}{lrr}
\toprule
\textbf{Metric} & \textbf{Value} & \textbf{Note} \\
\midrule
Accuracy (default threshold 0.5)     & 0.5455 & Binary decision used in main analysis \\
Best accuracy from threshold sweep   & 0.5909 & achieved at thresholds 0.35--0.65 \\
Area under the ROC curve (AUROC)     & 0.4132 & averaged across the three models \\
\bottomrule
\end{tabular}

\end{table}
\clearpage

\subsection*{Cellpose}
The following two tables present the per-image CPE probability scores and performance metrics obtained with CellPose.

\begin{table}[htbp]
\centering
\small
\setlength{\tabcolsep}{8pt}
\caption{CellPose automated CPE probability results}
\label{tab:cellpose-results}

\begin{tabular}{ccccc}
\toprule
\multicolumn{2}{c}{\textbf{Image}}
& \textbf{CRO}
& \multicolumn{2}{c}{\textbf{CellPose}} \\
\cmidrule(lr){1-2}\cmidrule(lr){3-3}\cmidrule(lr){4-5}
path & id
& CPE
& CPE Probability & CPE Detection \\
\midrule
\arrayrulecolor{gray!30}
1 & 101 & - & 0.6506 & 1 \\
\hline
1 & 103 & - & 0.6380 & 1 \\
\hline
1 & 104 & - & 0.7555 & 1 \\
\hline
1 & 201 & - & 0.6453 & 1 \\
\hline
1 & 202 & - & 0.7085 & 1 \\
\hline
1 & 305 & - & 0.6751 & 1 \\
\hline
1 & 401 & - & 0.6545 & 1 \\
\hline
1 & 406 & x & 0.7204 & 1 \\
\hline
1 & 501 & - & 0.6509 & 1 \\
\hline
2 & 101 & x & 0.6030 & 1 \\
\hline
2 & 201 & - & 0.6418 & 1 \\
\hline
2 & 202 & x & 0.6799 & 1 \\
\hline
2 & 303 & x & 0.6541 & 1 \\
\hline
2 & 307 & - & 0.6348 & 1 \\
\hline
2 & 308 & - & 0.7177 & 1 \\
\hline
2 & 402 & x & 0.6922 & 1 \\
\hline
2 & 405 & x & 0.6718 & 1 \\
\hline
2 & 406 & x & 0.7316 & 1 \\
\hline
2 & 503 & x & 0.6681 & 1 \\
\hline
2 & 504 & x & 0.7319 & 1 \\
\hline
2 & 507 & x & 0.6630 & 1 \\
\hline
2 & 508 & x & 0.6301 & 1 \\
\arrayrulecolor{black}
\bottomrule
\end{tabular}

\end{table}
\begin{table}[htbp]
\centering
\small
\caption{CellPose CPE detection performance metrics (description set, n=22).}
\label{tab:cellpose-performance}

\begin{tabular}{lrr}
\toprule
\textbf{Metric} & \textbf{Value} & \textbf{Note} \\
\midrule
Accuracy (default threshold 0.5) & 0.500 & Binary decision used in main analysis \\
Best accuracy from threshold sweep & 0.5909 & achieved at threshold = 0.65 \\
Area under the ROC curve (AUROC) & 0.6033 & overall discriminative ability \\
\bottomrule
\end{tabular}

\end{table}
\clearpage

\subsection*{Multimodal AI models}
The per-morphology CPE detection results for the four multimodal large language models (ChatGPT, Claude, Gemini, and Grok) are shown in Table \ref{tab:cpe_confusion}.

\clearpage
\begin{landscape}
\begin{table}[p]
\centering
\small
\setlength{\tabcolsep}{2pt}
\caption{LLM CPE detection results (with accessibility markers overlaid on colors)}
\label{tab:cpe_confusion}

\vspace{10pt}

\begin{tabular}{cc*{30}{>{\centering\arraybackslash}p{0.58cm}}}
\toprule
\multicolumn{2}{c}{\textbf{Image}} 
& \multicolumn{6}{c}{\textbf{CRO}} 
& \multicolumn{6}{c}{\textbf{ChatGPT}} 
& \multicolumn{6}{c}{\textbf{Claude}} 
& \multicolumn{6}{c}{\textbf{Gemini}} 
& \multicolumn{6}{c}{\textbf{Grok}} \\
\cmidrule(lr){1-2}\cmidrule(lr){3-8}\cmidrule(lr){9-14}\cmidrule(lr){15-20}\cmidrule(lr){21-26}\cmidrule(lr){27-32}
path & id 
& Dy & Ro & V & D & G & Re 
& Dy & Ro & V & D & G & Re 
& Dy & Ro & V & D & G & Re 
& Dy & Ro & V & D & G & Re 
& Dy & Ro & V & D & G & Re \\
\midrule
\arrayrulecolor{gray!30}
1 & 101 & \cellcolor{TNwhite}$\circ$ & \cellcolor{TNwhite}$\circ$ & \cellcolor{TNwhite}$\circ$ & \cellcolor{TNwhite}$\circ$ & \cellcolor{TNwhite}$\circ$ & \cellcolor{TNwhite}$\circ$
& \cellcolor{TNwhite}$\circ$ & \cellcolor{FPred}$\times$ & \cellcolor{FPred}$\times$ & \cellcolor{FPred}$\times$ & \cellcolor{TNwhite}$\circ$ & \cellcolor{FPred}$\times$
& \cellcolor{TNwhite}$\circ$ & \cellcolor{FPred}$\times$ & \cellcolor{TNwhite}$\circ$ & \cellcolor{FPred}$\times$ & \cellcolor{FPred}$\times$ & \cellcolor{FPred}$\times$
& \cellcolor{TNwhite}$\circ$ & \cellcolor{TNwhite}$\circ$ & \cellcolor{TNwhite}$\circ$ & \cellcolor{TNwhite}$\circ$ & \cellcolor{TNwhite}$\circ$ & \cellcolor{TNwhite}$\circ$
& \cellcolor{FPred}$\times$ & \cellcolor{FPred}$\times$ & \cellcolor{FPred}$\times$ & \cellcolor{FPred}$\times$ & \cellcolor{TNwhite}$\circ$ & \cellcolor{FPred}$\times$ \\
\hline
1 & 103 & \cellcolor{TNwhite}$\circ$ & \cellcolor{TNwhite}$\circ$ & \cellcolor{TNwhite}$\circ$ & \cellcolor{TNwhite}$\circ$ & \cellcolor{TNwhite}$\circ$ & \cellcolor{TNwhite}$\circ$
& \cellcolor{FPred}$\times$ & \cellcolor{FPred}$\times$ & \cellcolor{FPred}$\times$ & \cellcolor{FPred}$\times$ & \cellcolor{TNwhite}$\circ$ & \cellcolor{FPred}$\times$
& \cellcolor{TNwhite}$\circ$ & \cellcolor{FPred}$\times$ & \cellcolor{TNwhite}$\circ$ & \cellcolor{FPred}$\times$ & \cellcolor{FPred}$\times$ & \cellcolor{FPred}$\times$
& \cellcolor{TNwhite}$\circ$ & \cellcolor{TNwhite}$\circ$ & \cellcolor{TNwhite}$\circ$ & \cellcolor{TNwhite}$\circ$ & \cellcolor{TNwhite}$\circ$ & \cellcolor{TNwhite}$\circ$
& \cellcolor{FPred}$\times$ & \cellcolor{FPred}$\times$ & \cellcolor{FPred}$\times$ & \cellcolor{FPred}$\times$ & \cellcolor{TNwhite}$\circ$ & \cellcolor{FPred}$\times$ \\
\hline
1 & 104 & \cellcolor{TNwhite}$\circ$ & \cellcolor{TNwhite}$\circ$ & \cellcolor{TNwhite}$\circ$ & \cellcolor{TNwhite}$\circ$ & \cellcolor{TNwhite}$\circ$ & \cellcolor{TNwhite}$\circ$
& \cellcolor{FPred}$\times$ & \cellcolor{FPred}$\times$ & \cellcolor{FPred}$\times$ & \cellcolor{FPred}$\times$ & \cellcolor{FPred}$\times$ & \cellcolor{FPred}$\times$
& \cellcolor{TNwhite}$\circ$ & \cellcolor{FPred}$\times$ & \cellcolor{FPred}$\times$ & \cellcolor{FPred}$\times$ & \cellcolor{FPred}$\times$ & \cellcolor{FPred}$\times$
& \cellcolor{FPred}$\times$ & \cellcolor{FPred}$\times$ & \cellcolor{TNwhite}$\circ$ & \cellcolor{FPred}$\times$ & \cellcolor{FPred}$\times$ & \cellcolor{FPred}$\times$
& \cellcolor{FPred}$\times$ & \cellcolor{FPred}$\times$ & \cellcolor{FPred}$\times$ & \cellcolor{FPred}$\times$ & \cellcolor{FPred}$\times$ & \cellcolor{FPred}$\times$ \\
\hline
1 & 201 & \cellcolor{TNwhite}$\circ$ & \cellcolor{TNwhite}$\circ$ & \cellcolor{TNwhite}$\circ$ & \cellcolor{TNwhite}$\circ$ & \cellcolor{TNwhite}$\circ$ & \cellcolor{TNwhite}$\circ$
& \cellcolor{TNwhite}$\circ$ & \cellcolor{FPred}$\times$ & \cellcolor{FPred}$\times$ & \cellcolor{FPred}$\times$ & \cellcolor{TNwhite}$\circ$ & \cellcolor{FPred}$\times$
& \cellcolor{TNwhite}$\circ$ & \cellcolor{FPred}$\times$ & \cellcolor{TNwhite}$\circ$ & \cellcolor{FPred}$\times$ & \cellcolor{FPred}$\times$ & \cellcolor{FPred}$\times$
& \cellcolor{TNwhite}$\circ$ & \cellcolor{TNwhite}$\circ$ & \cellcolor{TNwhite}$\circ$ & \cellcolor{TNwhite}$\circ$ & \cellcolor{TNwhite}$\circ$ & \cellcolor{TNwhite}$\circ$
& \cellcolor{TNwhite}$\circ$ & \cellcolor{FPred}$\times$ & \cellcolor{FPred}$\times$ & \cellcolor{FPred}$\times$ & \cellcolor{TNwhite}$\circ$ & \cellcolor{FPred}$\times$ \\
\hline
1 & 202 & \cellcolor{TNwhite}$\circ$ & \cellcolor{TNwhite}$\circ$ & \cellcolor{TNwhite}$\circ$ & \cellcolor{TNwhite}$\circ$ & \cellcolor{TNwhite}$\circ$ & \cellcolor{TNwhite}$\circ$
& \cellcolor{FPred}$\times$ & \cellcolor{FPred}$\times$ & \cellcolor{FPred}$\times$ & \cellcolor{FPred}$\times$ & \cellcolor{FPred}$\times$ & \cellcolor{FPred}$\times$
& \cellcolor{TNwhite}$\circ$ & \cellcolor{FPred}$\times$ & \cellcolor{FPred}$\times$ & \cellcolor{FPred}$\times$ & \cellcolor{FPred}$\times$ & \cellcolor{FPred}$\times$
& \cellcolor{TNwhite}$\circ$ & \cellcolor{FPred}$\times$ & \cellcolor{TNwhite}$\circ$ & \cellcolor{TNwhite}$\circ$ & \cellcolor{TNwhite}$\circ$ & \cellcolor{FPred}$\times$
& \cellcolor{FPred}$\times$ & \cellcolor{FPred}$\times$ & \cellcolor{FPred}$\times$ & \cellcolor{FPred}$\times$ & \cellcolor{FPred}$\times$ & \cellcolor{FPred}$\times$ \\
\hline
1 & 305 & \cellcolor{TNwhite}$\circ$ & \cellcolor{TNwhite}$\circ$ & \cellcolor{TNwhite}$\circ$ & \cellcolor{TNwhite}$\circ$ & \cellcolor{TNwhite}$\circ$ & \cellcolor{TNwhite}$\circ$
& \cellcolor{TNwhite}$\circ$ & \cellcolor{FPred}$\times$ & \cellcolor{FPred}$\times$ & \cellcolor{TNwhite}$\circ$ & \cellcolor{TNwhite}$\circ$ & \cellcolor{FPred}$\times$
& \cellcolor{TNwhite}$\circ$ & \cellcolor{FPred}$\times$ & \cellcolor{FPred}$\times$ & \cellcolor{FPred}$\times$ & \cellcolor{FPred}$\times$ & \cellcolor{FPred}$\times$
& \cellcolor{TNwhite}$\circ$ & \cellcolor{TNwhite}$\circ$ & \cellcolor{TNwhite}$\circ$ & \cellcolor{TNwhite}$\circ$ & \cellcolor{TNwhite}$\circ$ & \cellcolor{TNwhite}$\circ$
& \cellcolor{TNwhite}$\circ$ & \cellcolor{FPred}$\times$ & \cellcolor{FPred}$\times$ & \cellcolor{FPred}$\times$ & \cellcolor{TNwhite}$\circ$ & \cellcolor{FPred}$\times$ \\
\hline
1 & 401 & \cellcolor{TNwhite}$\circ$ & \cellcolor{TNwhite}$\circ$ & \cellcolor{TNwhite}$\circ$ & \cellcolor{TNwhite}$\circ$ & \cellcolor{TNwhite}$\circ$ & \cellcolor{TNwhite}$\circ$
& \cellcolor{TNwhite}$\circ$ & \cellcolor{FPred}$\times$ & \cellcolor{TNwhite}$\circ$ & \cellcolor{TNwhite}$\circ$ & \cellcolor{TNwhite}$\circ$ & \cellcolor{FPred}$\times$
& \cellcolor{TNwhite}$\circ$ & \cellcolor{FPred}$\times$ & \cellcolor{TNwhite}$\circ$ & \cellcolor{FPred}$\times$ & \cellcolor{FPred}$\times$ & \cellcolor{FPred}$\times$
& \cellcolor{TNwhite}$\circ$ & \cellcolor{TNwhite}$\circ$ & \cellcolor{TNwhite}$\circ$ & \cellcolor{TNwhite}$\circ$ & \cellcolor{TNwhite}$\circ$ & \cellcolor{TNwhite}$\circ$
& \cellcolor{TNwhite}$\circ$ & \cellcolor{FPred}$\times$ & \cellcolor{FPred}$\times$ & \cellcolor{TNwhite}$\circ$ & \cellcolor{TNwhite}$\circ$ & \cellcolor{FPred}$\times$ \\
\hline
1 & 406 & \cellcolor{TNwhite}$\circ$ & \cellcolor{TNwhite}$\circ$ & \cellcolor{TPgreen}$+$ & \cellcolor{TNwhite}$\circ$ & \cellcolor{TNwhite}$\circ$ & \cellcolor{TNwhite}$\circ$
& \cellcolor{TNwhite}$\circ$ & \cellcolor{TNwhite}$\circ$ & \cellcolor{FNyellow}$-$ & \cellcolor{TNwhite}$\circ$ & \cellcolor{FPred}$\times$ & \cellcolor{FPred}$\times$
& \cellcolor{TNwhite}$\circ$ & \cellcolor{FPred}$\times$ & \cellcolor{TPgreen}$+$ & \cellcolor{FPred}$\times$ & \cellcolor{FPred}$\times$ & \cellcolor{FPred}$\times$
& \cellcolor{TNwhite}$\circ$ & \cellcolor{TNwhite}$\circ$ & \cellcolor{FNyellow}$-$ & \cellcolor{TNwhite}$\circ$ & \cellcolor{TNwhite}$\circ$ & \cellcolor{TNwhite}$\circ$
& \cellcolor{TNwhite}$\circ$ & \cellcolor{TNwhite}$\circ$ & \cellcolor{TPgreen}$+$ & \cellcolor{TNwhite}$\circ$ & \cellcolor{TNwhite}$\circ$ & \cellcolor{FPred}$\times$ \\
\hline
1 & 501 & \cellcolor{TNwhite}$\circ$ & \cellcolor{TNwhite}$\circ$ & \cellcolor{TNwhite}$\circ$ & \cellcolor{TNwhite}$\circ$ & \cellcolor{TNwhite}$\circ$ & \cellcolor{TNwhite}$\circ$
& \cellcolor{TNwhite}$\circ$ & \cellcolor{FPred}$\times$ & \cellcolor{TNwhite}$\circ$ & \cellcolor{TNwhite}$\circ$ & \cellcolor{TNwhite}$\circ$ & \cellcolor{FPred}$\times$
& \cellcolor{TNwhite}$\circ$ & \cellcolor{FPred}$\times$ & \cellcolor{TNwhite}$\circ$ & \cellcolor{FPred}$\times$ & \cellcolor{FPred}$\times$ & \cellcolor{FPred}$\times$
& \cellcolor{TNwhite}$\circ$ & \cellcolor{TNwhite}$\circ$ & \cellcolor{TNwhite}$\circ$ & \cellcolor{TNwhite}$\circ$ & \cellcolor{TNwhite}$\circ$ & \cellcolor{TNwhite}$\circ$
& \cellcolor{TNwhite}$\circ$ & \cellcolor{FPred}$\times$ & \cellcolor{TNwhite}$\circ$ & \cellcolor{TNwhite}$\circ$ & \cellcolor{TNwhite}$\circ$ & \cellcolor{FPred}$\times$ \\
\hline
2 & 101 & \cellcolor{TNwhite}$\circ$ & \cellcolor{TPgreen}$+$ & \cellcolor{TNwhite}$\circ$ & \cellcolor{TNwhite}$\circ$ & \cellcolor{TNwhite}$\circ$ & \cellcolor{TNwhite}$\circ$
& \cellcolor{FPred}$\times$ & \cellcolor{TPgreen}$+$ & \cellcolor{FPred}$\times$ & \cellcolor{FPred}$\times$ & \cellcolor{TNwhite}$\circ$ & \cellcolor{FPred}$\times$
& \cellcolor{TNwhite}$\circ$ & \cellcolor{TPgreen}$+$ & \cellcolor{TNwhite}$\circ$ & \cellcolor{FPred}$\times$ & \cellcolor{FPred}$\times$ & \cellcolor{FPred}$\times$
& \cellcolor{TNwhite}$\circ$ & \cellcolor{FNyellow}$-$ & \cellcolor{TNwhite}$\circ$ & \cellcolor{TNwhite}$\circ$ & \cellcolor{TNwhite}$\circ$ & \cellcolor{TNwhite}$\circ$
& \cellcolor{TNwhite}$\circ$ & \cellcolor{TPgreen}$+$ & \cellcolor{FPred}$\times$ & \cellcolor{FPred}$\times$ & \cellcolor{TNwhite}$\circ$ & \cellcolor{FPred}$\times$ \\
\hline
2 & 201 & \cellcolor{TNwhite}$\circ$ & \cellcolor{TNwhite}$\circ$ & \cellcolor{TNwhite}$\circ$ & \cellcolor{TNwhite}$\circ$ & \cellcolor{TNwhite}$\circ$ & \cellcolor{TNwhite}$\circ$
& \cellcolor{FPred}$\times$ & \cellcolor{FPred}$\times$ & \cellcolor{FPred}$\times$ & \cellcolor{FPred}$\times$ & \cellcolor{TNwhite}$\circ$ & \cellcolor{FPred}$\times$
& \cellcolor{TNwhite}$\circ$ & \cellcolor{FPred}$\times$ & \cellcolor{FPred}$\times$ & \cellcolor{FPred}$\times$ & \cellcolor{FPred}$\times$ & \cellcolor{FPred}$\times$
& \cellcolor{TNwhite}$\circ$ & \cellcolor{TNwhite}$\circ$ & \cellcolor{TNwhite}$\circ$ & \cellcolor{TNwhite}$\circ$ & \cellcolor{TNwhite}$\circ$ & \cellcolor{TNwhite}$\circ$
& \cellcolor{FPred}$\times$ & \cellcolor{FPred}$\times$ & \cellcolor{FPred}$\times$ & \cellcolor{FPred}$\times$ & \cellcolor{TNwhite}$\circ$ & \cellcolor{FPred}$\times$ \\
\hline
2 & 202 & \cellcolor{TNwhite}$\circ$ & \cellcolor{TPgreen}$+$ & \cellcolor{TNwhite}$\circ$ & \cellcolor{TNwhite}$\circ$ & \cellcolor{TNwhite}$\circ$ & \cellcolor{TPgreen}$+$
& \cellcolor{FPred}$\times$ & \cellcolor{TPgreen}$+$ & \cellcolor{FPred}$\times$ & \cellcolor{FPred}$\times$ & \cellcolor{TNwhite}$\circ$ & \cellcolor{TPgreen}$+$
& \cellcolor{TNwhite}$\circ$ & \cellcolor{TPgreen}$+$ & \cellcolor{TNwhite}$\circ$ & \cellcolor{FPred}$\times$ & \cellcolor{FPred}$\times$ & \cellcolor{TPgreen}$+$
& \cellcolor{FPred}$\times$ & \cellcolor{TPgreen}$+$ & \cellcolor{TNwhite}$\circ$ & \cellcolor{FPred}$\times$ & \cellcolor{TNwhite}$\circ$ & \cellcolor{TPgreen}$+$
& \cellcolor{FPred}$\times$ & \cellcolor{TPgreen}$+$ & \cellcolor{FPred}$\times$ & \cellcolor{FPred}$\times$ & \cellcolor{TNwhite}$\circ$ & \cellcolor{TPgreen}$+$ \\
\hline
2 & 303 & \cellcolor{TNwhite}$\circ$ & \cellcolor{TPgreen}$+$ & \cellcolor{TPgreen}$+$ & \cellcolor{TNwhite}$\circ$ & \cellcolor{TNwhite}$\circ$ & \cellcolor{TNwhite}$\circ$
& \cellcolor{FPred}$\times$ & \cellcolor{TPgreen}$+$ & \cellcolor{TPgreen}$+$ & \cellcolor{FPred}$\times$ & \cellcolor{TNwhite}$\circ$ & \cellcolor{FPred}$\times$
& \cellcolor{FPred}$\times$ & \cellcolor{TPgreen}$+$ & \cellcolor{FNyellow}$-$ & \cellcolor{FPred}$\times$ & \cellcolor{FPred}$\times$ & \cellcolor{FPred}$\times$
& \cellcolor{TNwhite}$\circ$ & \cellcolor{FNyellow}$-$ & \cellcolor{FNyellow}$-$ & \cellcolor{TNwhite}$\circ$ & \cellcolor{TNwhite}$\circ$ & \cellcolor{TNwhite}$\circ$
& \cellcolor{FPred}$\times$ & \cellcolor{TPgreen}$+$ & \cellcolor{TPgreen}$+$ & \cellcolor{FPred}$\times$ & \cellcolor{TNwhite}$\circ$ & \cellcolor{FPred}$\times$ \\
\hline
2 & 307 & \cellcolor{TNwhite}$\circ$ & \cellcolor{TNwhite}$\circ$ & \cellcolor{TNwhite}$\circ$ & \cellcolor{TNwhite}$\circ$ & \cellcolor{TNwhite}$\circ$ & \cellcolor{TNwhite}$\circ$
& \cellcolor{FPred}$\times$ & \cellcolor{FPred}$\times$ & \cellcolor{FPred}$\times$ & \cellcolor{FPred}$\times$ & \cellcolor{TNwhite}$\circ$ & \cellcolor{FPred}$\times$
& \cellcolor{TNwhite}$\circ$ & \cellcolor{FPred}$\times$ & \cellcolor{TNwhite}$\circ$ & \cellcolor{FPred}$\times$ & \cellcolor{FPred}$\times$ & \cellcolor{FPred}$\times$
& \cellcolor{FPred}$\times$ & \cellcolor{FPred}$\times$ & \cellcolor{TNwhite}$\circ$ & \cellcolor{TNwhite}$\circ$ & \cellcolor{FPred}$\times$ & \cellcolor{FPred}$\times$
& \cellcolor{FPred}$\times$ & \cellcolor{FPred}$\times$ & \cellcolor{TNwhite}$\circ$ & \cellcolor{FPred}$\times$ & \cellcolor{TNwhite}$\circ$ & \cellcolor{FPred}$\times$ \\
\hline
2 & 308 & \cellcolor{TNwhite}$\circ$ & \cellcolor{TNwhite}$\circ$ & \cellcolor{TNwhite}$\circ$ & \cellcolor{TNwhite}$\circ$ & \cellcolor{TNwhite}$\circ$ & \cellcolor{TNwhite}$\circ$
& \cellcolor{FPred}$\times$ & \cellcolor{FPred}$\times$ & \cellcolor{TNwhite}$\circ$ & \cellcolor{FPred}$\times$ & \cellcolor{TNwhite}$\circ$ & \cellcolor{FPred}$\times$
& \cellcolor{TNwhite}$\circ$ & \cellcolor{FPred}$\times$ & \cellcolor{TNwhite}$\circ$ & \cellcolor{FPred}$\times$ & \cellcolor{FPred}$\times$ & \cellcolor{FPred}$\times$
& \cellcolor{FPred}$\times$ & \cellcolor{FPred}$\times$ & \cellcolor{FPred}$\times$ & \cellcolor{TNwhite}$\circ$ & \cellcolor{FPred}$\times$ & \cellcolor{FPred}$\times$
& \cellcolor{FPred}$\times$ & \cellcolor{FPred}$\times$ & \cellcolor{TNwhite}$\circ$ & \cellcolor{FPred}$\times$ & \cellcolor{TNwhite}$\circ$ & \cellcolor{FPred}$\times$ \\
\hline
2 & 402 & \cellcolor{TPgreen}$+$ & \cellcolor{TPgreen}$+$ & \cellcolor{TPgreen}$+$ & \cellcolor{TNwhite}$\circ$ & \cellcolor{TNwhite}$\circ$ & \cellcolor{TPgreen}$+$
& \cellcolor{TPgreen}$+$ & \cellcolor{TPgreen}$+$ & \cellcolor{TPgreen}$+$ & \cellcolor{FPred}$\times$ & \cellcolor{TNwhite}$\circ$ & \cellcolor{TPgreen}$+$
& \cellcolor{FNyellow}$-$ & \cellcolor{TPgreen}$+$ & \cellcolor{TPgreen}$+$ & \cellcolor{FPred}$\times$ & \cellcolor{FPred}$\times$ & \cellcolor{TPgreen}$+$
& \cellcolor{FNyellow}$-$ & \cellcolor{FNyellow}$-$ & \cellcolor{FNyellow}$-$ & \cellcolor{TNwhite}$\circ$ & \cellcolor{TNwhite}$\circ$ & \cellcolor{FNyellow}$-$
& \cellcolor{TPgreen}$+$ & \cellcolor{TPgreen}$+$ & \cellcolor{TPgreen}$+$ & \cellcolor{FPred}$\times$ & \cellcolor{TNwhite}$\circ$ & \cellcolor{TPgreen}$+$ \\
\hline
2 & 405 & \cellcolor{TPgreen}$+$ & \cellcolor{TPgreen}$+$ & \cellcolor{TNwhite}$\circ$ & \cellcolor{TNwhite}$\circ$ & \cellcolor{TPgreen}$+$ & \cellcolor{TNwhite}$\circ$
& \cellcolor{TPgreen}$+$ & \cellcolor{TPgreen}$+$ & \cellcolor{FPred}$\times$ & \cellcolor{FPred}$\times$ & \cellcolor{FNyellow}$-$ & \cellcolor{FPred}$\times$
& \cellcolor{FNyellow}$-$ & \cellcolor{TPgreen}$+$ & \cellcolor{TNwhite}$\circ$ & \cellcolor{FPred}$\times$ & \cellcolor{TPgreen}$+$ & \cellcolor{FPred}$\times$
& \cellcolor{FNyellow}$-$ & \cellcolor{FNyellow}$-$ & \cellcolor{TNwhite}$\circ$ & \cellcolor{TNwhite}$\circ$ & \cellcolor{FNyellow}$-$ & \cellcolor{TNwhite}$\circ$
& \cellcolor{TPgreen}$+$ & \cellcolor{TPgreen}$+$ & \cellcolor{FPred}$\times$ & \cellcolor{FPred}$\times$ & \cellcolor{FNyellow}$-$ & \cellcolor{FPred}$\times$ \\
\hline
2 & 406 & \cellcolor{TPgreen}$+$ & \cellcolor{TPgreen}$+$ & \cellcolor{TNwhite}$\circ$ & \cellcolor{TNwhite}$\circ$ & \cellcolor{TPgreen}$+$ & \cellcolor{TNwhite}$\circ$
& \cellcolor{TPgreen}$+$ & \cellcolor{TPgreen}$+$ & \cellcolor{FPred}$\times$ & \cellcolor{FPred}$\times$ & \cellcolor{FNyellow}$-$ & \cellcolor{FPred}$\times$
& \cellcolor{FNyellow}$-$ & \cellcolor{TPgreen}$+$ & \cellcolor{FPred}$\times$ & \cellcolor{FPred}$\times$ & \cellcolor{TPgreen}$+$ & \cellcolor{FPred}$\times$
& \cellcolor{FNyellow}$-$ & \cellcolor{FNyellow}$-$ & \cellcolor{TNwhite}$\circ$ & \cellcolor{TNwhite}$\circ$ & \cellcolor{FNyellow}$-$ & \cellcolor{TNwhite}$\circ$
& \cellcolor{TPgreen}$+$ & \cellcolor{TPgreen}$+$ & \cellcolor{FPred}$\times$ & \cellcolor{FPred}$\times$ & \cellcolor{FNyellow}$-$ & \cellcolor{FPred}$\times$ \\
\hline
2 & 503 & \cellcolor{TPgreen}$+$ & \cellcolor{TNwhite}$\circ$ & \cellcolor{TPgreen}$+$ & \cellcolor{TPgreen}$+$ & \cellcolor{TPgreen}$+$ & \cellcolor{TPgreen}$+$
& \cellcolor{TPgreen}$+$ & \cellcolor{FPred}$\times$ & \cellcolor{TPgreen}$+$ & \cellcolor{TPgreen}$+$ & \cellcolor{FNyellow}$-$ & \cellcolor{TPgreen}$+$
& \cellcolor{FNyellow}$-$ & \cellcolor{FPred}$\times$ & \cellcolor{TPgreen}$+$ & \cellcolor{TPgreen}$+$ & \cellcolor{TPgreen}$+$ & \cellcolor{TPgreen}$+$
& \cellcolor{FNyellow}$-$ & \cellcolor{TNwhite}$\circ$ & \cellcolor{FNyellow}$-$ & \cellcolor{FNyellow}$-$ & \cellcolor{FNyellow}$-$ & \cellcolor{FNyellow}$-$
& \cellcolor{TPgreen}$+$ & \cellcolor{FPred}$\times$ & \cellcolor{TPgreen}$+$ & \cellcolor{TPgreen}$+$ & \cellcolor{FNyellow}$-$ & \cellcolor{TPgreen}$+$ \\
\hline
2 & 504 & \cellcolor{TPgreen}$+$ & \cellcolor{TNwhite}$\circ$ & \cellcolor{TPgreen}$+$ & \cellcolor{TPgreen}$+$ & \cellcolor{TPgreen}$+$ & \cellcolor{TPgreen}$+$
& \cellcolor{TPgreen}$+$ & \cellcolor{FPred}$\times$ & \cellcolor{TPgreen}$+$ & \cellcolor{FNyellow}$-$ & \cellcolor{FNyellow}$-$ & \cellcolor{TPgreen}$+$
& \cellcolor{FNyellow}$-$ & \cellcolor{FPred}$\times$ & \cellcolor{TPgreen}$+$ & \cellcolor{TPgreen}$+$ & \cellcolor{TPgreen}$+$ & \cellcolor{TPgreen}$+$
& \cellcolor{FNyellow}$-$ & \cellcolor{TNwhite}$\circ$ & \cellcolor{FNyellow}$-$ & \cellcolor{FNyellow}$-$ & \cellcolor{FNyellow}$-$ & \cellcolor{FNyellow}$-$
& \cellcolor{TPgreen}$+$ & \cellcolor{FPred}$\times$ & \cellcolor{TPgreen}$+$ & \cellcolor{FNyellow}$-$ & \cellcolor{FNyellow}$-$ & \cellcolor{TPgreen}$+$ \\
\hline
2 & 507 & \cellcolor{TPgreen}$+$ & \cellcolor{TPgreen}$+$ & \cellcolor{TNwhite}$\circ$ & \cellcolor{TNwhite}$\circ$ & \cellcolor{TPgreen}$+$ & \cellcolor{TNwhite}$\circ$
& \cellcolor{TPgreen}$+$ & \cellcolor{TPgreen}$+$ & \cellcolor{FPred}$\times$ & \cellcolor{FPred}$\times$ & \cellcolor{FNyellow}$-$ & \cellcolor{FPred}$\times$
& \cellcolor{FNyellow}$-$ & \cellcolor{TPgreen}$+$ & \cellcolor{FPred}$\times$ & \cellcolor{FPred}$\times$ & \cellcolor{TPgreen}$+$ & \cellcolor{FPred}$\times$
& \cellcolor{FNyellow}$-$ & \cellcolor{FNyellow}$-$ & \cellcolor{TNwhite}$\circ$ & \cellcolor{TNwhite}$\circ$ & \cellcolor{FNyellow}$-$ & \cellcolor{TNwhite}$\circ$
& \cellcolor{FNyellow}$-$ & \cellcolor{TPgreen}$+$ & \cellcolor{FPred}$\times$ & \cellcolor{FPred}$\times$ & \cellcolor{FNyellow}$-$ & \cellcolor{FPred}$\times$ \\
\hline
2 & 508 & \cellcolor{TPgreen}$+$ & \cellcolor{TPgreen}$+$ & \cellcolor{TNwhite}$\circ$ & \cellcolor{TNwhite}$\circ$ & \cellcolor{TPgreen}$+$ & \cellcolor{TNwhite}$\circ$
& \cellcolor{TPgreen}$+$ & \cellcolor{TPgreen}$+$ & \cellcolor{FPred}$\times$ & \cellcolor{FPred}$\times$ & \cellcolor{FNyellow}$-$ & \cellcolor{FPred}$\times$
& \cellcolor{FNyellow}$-$ & \cellcolor{TPgreen}$+$ & \cellcolor{FPred}$\times$ & \cellcolor{FPred}$\times$ & \cellcolor{TPgreen}$+$ & \cellcolor{FPred}$\times$
& \cellcolor{FNyellow}$-$ & \cellcolor{FNyellow}$-$ & \cellcolor{TNwhite}$\circ$ & \cellcolor{TNwhite}$\circ$ & \cellcolor{FNyellow}$-$ & \cellcolor{TNwhite}$\circ$
& \cellcolor{TPgreen}$+$ & \cellcolor{TPgreen}$+$ & \cellcolor{FPred}$\times$ & \cellcolor{FPred}$\times$ & \cellcolor{FNyellow}$-$ & \cellcolor{FPred}$\times$ \\
\arrayrulecolor{black}
\bottomrule
\end{tabular}
\end{table}

% ====================== HORIZONTAL LEGEND ======================
\begin{table}[p]
\centering
\begin{tabular}{>{\centering\arraybackslash}p{18pt} l >{\centering\arraybackslash}p{18pt} l >{\centering\arraybackslash}p{18pt} l >{\centering\arraybackslash}p{18pt} l}
\toprule
\cellcolor{TPgreen}$+$ & True Positive 
& \cellcolor{TNwhite}$\circ$ & True Negative 
& \cellcolor{FNyellow}$-$ & False Negative 
& \cellcolor{FPred}$\times$ & False Positive \\
\bottomrule
\end{tabular}
\end{table}

\end{landscape}
\clearpage

\subsection*{Aggregate performance summary}
Overall performance metrics of AIRVIC, DVICE, CellPose, and the four multimodal models are summarized in Table \ref{tab:aggregate-results}. AIRVIC, DVICE, and CellPose are compared against simple always-true and always-false binary classifiers. The multimodal models are compared against always-true and always-false CPE-type classifiers. The accuracy of the always-false CPE-type classifier is a property of the dataset; any useful model must substantially outperform this baseline.

\begin{table}[htbp]
\centering
\caption{Performance metrics of AIRVIC, CellPose, and four large language models for cytopathic effect (CPE) detection. For the large language models, false-positive rate (FP rate), false-negative rate (FN rate), true-positive rate (TP rate), true-negative rate (TN rate), and overall accuracy were first calculated separately for each of the six CPE morphological types and then macro-averaged. For AIRVIC and CellPose, metrics were computed directly on the binary CPE detection task (CellPose probabilities were binarized at the conventional threshold of 0.5).}
\label{tab:cpe-performance-metrics}

\begin{tabular}{lrrrrr}
\toprule
\textbf{Model} & \textbf{FP rate} & \textbf{FN rate} & \textbf{TP rate} & \textbf{TN rate} & \textbf{Overall Accuracy} \\
\midrule
\arrayrulecolor{gray!30}
AIRVIC     & 0.9091 & 0.0000 & 1.0000 & 0.0909 & 0.5455 \\
\hline
CellPose   & 1.0000 & 0.0000 & 1.0000 & 0.0000 & 0.5000 \\
\hline
ChatGPT    & 0.5520 & 0.3506 & 0.6494 & 0.4480 & 0.5588 \\
\hline
Claude     & 0.5902 & 0.3275 & 0.6725 & 0.4098 & 0.5425 \\
\hline
Gemini     & 0.1283 & 0.8468 & 0.1532 & 0.8717 & 0.7092 \\
\hline
Grok       & 0.5670 & 0.3566 & 0.6434 & 0.4330 & 0.5490 \\
\arrayrulecolor{black}
\bottomrule
\end{tabular}

\end{table}
\section*{Discussion}
\label{sec:discussion}
\subsection*{Domain-specific tools}
AIRVIC detected CPE in 21 of the 22 images in the description set and in 100 of the 101 total images, resulting in a high false-positive rate of 0.909 and an overall accuracy of only 0.545. The test images are visually similar to the Vero-cell images used in AIRVIC’s original training set, as demonstrated by table \ref{tab:airvic-image-comparison}. Despite this similarity, the model failed to generalize effectively to the present dataset. Although direct access to the AIRVIC training data was not available, it is possible that the generally low confluency of the current dataset was not well represented in the original training set.

\clearpage
\begin{landscape}
\begin{table}[p]
\centering
\small
\setlength{\tabcolsep}{7pt}
\renewcommand{\arraystretch}{1.55}
\caption{Comparison of image acquisition and processing parameters between the AIRVIC training dataset \cite{airvic} and the current Vero cell culture dataset \cite{mathews2025}.}
\label{tab:airvic-image-comparison}
\vspace{10pt}

\begin{tabular}{p{4.2cm} p{5.3cm} p{5.3cm} p{9.0cm}}
\toprule
\textbf{Aspect} & \textbf{AIRVIC dataset} & \textbf{Vero cell
culture dataset} & \textbf{Potential impact on AIRVIC accuracy} \\
\midrule
\arrayrulecolor{gray!30}
Microscope system & ZEISS Axio Observer 5 widefield microscope & EVOS\textsuperscript{\textregistered} FL Auto system (Thermo Fisher, Cat. no. AMAFD1000) & Moderate — model trained exclusively on a single microscope brand and optical path; cross-platform generalization has not been validated \\
\hline
Objective magnification & 5×, 10×, 20×, 40× (phase contrast) & 10× and 20× (AMG LPlan FL PH objectives) & Low — strong overlap with the optimal training magnifications \\
\hline
Contrast / Illumination & Phase contrast only (Ph1/Ph2 phase rings) & Phase contrast only & Low — excellent match with training conditions \\
\hline
Camera & ZEISS Axiocam 506 color CMOS (2752 × 2208 px) & High-sensitivity CMOS (EVOS) & Moderate — differences in sensor response, color balance, and pixel characteristics possible \\
\hline
Original file format & JPEG (24-bit color) & TIFF & Negligible \\
\hline
Final format submitted to AIRVIC & JPEG (internally converted to grayscale) & PNG (converted from TIFF, color preserved) & Low–Moderate — model expects grayscale-like input; absence of explicit grayscale conversion may introduce minor color-related bias \\
\hline
Acquisition software & ZEISS ZEN lite 2.3 & DIBImage v1.0.0 (rev. 462) / Evos software & Low \\
\arrayrulecolor{black}
\bottomrule
\end{tabular}

\end{table}
\end{landscape}
\clearpage

The three pre-trained DVICE models performed similarly poorly, yielding a high false-positive rate of 0.818 and an overall accuracy of only 0.545. When the average probability across the three models was used, DVICE achieved a maximum accuracy of 0.5909 (AUROC = 0.4132) even after threshold optimization. The models produced many false positives on CPE-negative images while missing some true positives. It should be noted that the DVICE training set consisted of transmitted-light microscope images, whereas the current dataset consists of phase-contrast images (see Table \ref{tab:dvice-image-comparison}). This difference in imaging modality alone may account for the poor performance of the DVICE models in this study. The low confluency of the current dataset is unlikely to explain the poor performance, since the DVICE authors explicitly state that their process enabled ``the network to learn that low cell confluency is not a defining hallmark of viral infection state'' \cite{dvice}.

\clearpage
\begin{landscape}
\begin{table}[p]
\centering
\small
\setlength{\tabcolsep}{7pt}
\renewcommand{\arraystretch}{1.55}
\caption{Comparison of image acquisition and processing parameters between the DVICE training dataset\cite{dvice} and the current Vero cell culture dataset\cite{mathews2025}.}
\label{tab:dvice-image-comparison}
\vspace{10pt}

\begin{tabular}{p{4.2cm} p{5.3cm} p{5.3cm} p{9.0cm}}
\toprule
\textbf{Aspect} & \textbf{DVICE expectation (training data)} & \textbf{Current dataset (this study)} & \textbf{Potential impact on DVICE accuracy} \\
\midrule
\arrayrulecolor{gray!30}
Microscope system & ImageXpress Micro Confocal (IXM-C, Molecular Devices); Cytation 5 (validation) & EVOS\textsuperscript{\textregistered} FL Auto system (Thermo Fisher, Cat. no. AMAFD1000) & Moderate — model trained on a different high-throughput confocal system; cross-platform generalization not validated \\
\hline
Objective magnification & 4× air objective & 10× and 20× (AMG LPlan FL PH objectives) & Moderate — lower magnification and different optical setup than training data \\
\hline
Contrast / Illumination & Transmitted light (TL), label-free & Phase contrast only & High — fundamental difference in contrast mechanism (brightfield TL vs phase contrast) \\
\hline
Camera / Resolution & 2048 × 2048 px (16-bit) or 1992 × 1992 px (Cytation 5) & High-sensitivity CMOS (EVOS) & Moderate — differences in sensor response, bit depth, and pixel characteristics \\
\hline
Original file format & 16-bit raw (proprietary) & TIFF & Negligible \\
\hline
Final format / preprocessing & Rescaled to 8-bit PNG + resized to 224 × 224 px (RGB) & PNG (converted from TIFF, color preserved) & Moderate — model expects specific 224 × 224 RGB input after min-max normalization; current pipeline does not match this preprocessing \\
\hline
Cell line & VeroE6 / VeroE6-TMPRSS2 (SARS-CoV-2) + multiple human lines & Vero (standard) & Low — cell line is comparable \\
\hline
Acquisition timing & 7 days post-infection (dpi) & Variable (multiple passages) & Moderate — CPE morphology may differ due to timing and passage number \\
\arrayrulecolor{black}
\bottomrule
\end{tabular}

\end{table}
\end{landscape}
\clearpage

Cellpose similarly classified nearly all images as CPE-positive (accuracy of 0.500 at the default threshold of 0.5, with a maximum accuracy of 0.5909). Consequently, the morphological-proxy approach did not provide useful discrimination in this dataset.

\subsection*{Multimodal AI tools}
The four general-purpose multimodal models exhibited macro-averaged accuracies between 0.542 and 0.709. Gemini achieved the highest overall accuracy (0.709) but at the cost of a very high false-negative rate, effectively behaving like a conservative ``no-CPE'' classifier; this performance was still lower than that of the simple always-negative baseline model. ChatGPT, Claude, and Grok showed poor performance (accuracies 0.549--0.559) with substantial rates of both false positives and false negatives. None of the multimodal models reached a level of reliability that would support their use as objective CPE detection tools for this dataset.

\subsection*{Limitations}
The ground-truth labels are based on narrative descriptions provided by the CRO that were not generated under a CPE-specific protocol. The description set is small (\(n = 22\)), and the remaining 83 images could not be evaluated quantitatively. Domain-specific tools were tested exactly as publicly available, with no retraining or fine-tuning performed. Prompt engineering for the multimodal models, although extensive, remains inherently model-dependent.

Furthermore, the images from the experimental stage exhibited relatively low cell confluency compared with typical culture conditions under which CPE is normally observed. It is possible that the training sets of the domain-specific models did not include images with similarly low confluency.

\subsection*{Implications}
None of the evaluated AI tools demonstrated sufficient accuracy or robustness to replace or act as an objective second opinion for human expert interpretation of CPE in Vero-cell cultures under the conditions tested. As a consequence, an objective assessment of CPE presence could not be obtained for the remaining 83 images of the EXP stage. Future work may benefit from larger, prospectively labeled datasets and from domain-specific fine-tuning of multimodal models.

\nolinenumbers

\bibliography{references}

\end{document}

